\documentclass[12pt]{article}
\usepackage[utf8]{inputenc}
\usepackage{amsmath, amssymb, amsfonts}
\usepackage{geometry}
\usepackage{booktabs}
\usepackage{enumitem}
\geometry{margin=1in}
\setlist[enumerate]{leftmargin=*}

\title{Practice Problem Set 5\\ \vspace{15pt} \large ECON 101 -- Macroeconomics}
\author{}
\date{Fall 2025}

\begin{document}
\maketitle



\section*{Part A – Multiple Choice Questions}

\textbf{1.}  
A decrease in the Canadian price level, holding everything else constant, will:  \\

(A) Shift the aggregate demand curve to the left.  
(B) Shift the aggregate demand curve to the right.  
(C) Cause a movement \emph{down} along the aggregate demand curve.  
(D) Cause a movement \emph{up} along the aggregate demand curve.\\



\textbf{2.}  
Which of the following will \emph{increase} aggregate demand in Canada?  \\

(A) A fall in consumer confidence.  
(B) A rise in the Bank of Canada’s policy interest rate.  
(C) A depreciation of the Canadian dollar.  
(D) A decrease in government purchases.\\



\textbf{3.}  
In the short run, an increase in the price level will:  \\

(A) Increase firms’ real profits and shift SRAS to the right.  
(B) Cause a movement \emph{along} the short-run aggregate supply curve.  
(C) Shift the SRAS curve to the left.  
(D) Have no effect on output supplied.\\



\textbf{4.}  
If the required reserve ratio is 10\% and a new deposit of \$10,000 is made in a commercial bank, the \emph{maximum} possible increase in the money supply is:  \\

(A) \$1,000 \quad (B) \$9,000 \quad (C) \$100,000 \quad (D) \$90,000\\



\textbf{5.}  
When the Bank of Canada buys government securities on the open market, it is:  \\

(A) Increasing the money supply and lowering interest rates.  
(B) Decreasing the money supply and raising interest rates.  
(C) Reducing bank reserves and restricting credit.  
(D) Attempting to shift aggregate supply.

\newpage

\section*{Part B – Long Analytical Problem}

\textbf{6. Monetary Policy and Aggregate Demand–Aggregate Supply Interaction}

Suppose the economy of Canada is initially in short-run equilibrium where potential GDP (\(Y^*\)) equals \$2{,}000 billion and the price level (\(P\)) equals 100.  

An unexpected fall in consumer confidence causes autonomous consumption to decline by \$100 billion, shifting the aggregate demand (AD) curve leftward.

\textbf{Tasks:}
\begin{enumerate}
    \item[(a)] Explain qualitatively what happens to real GDP and the price level in the short run. Illustrate using the AD–AS model.
    \item[(b)] If the marginal propensity to consume (MPC) is 0.8, calculate the total change in equilibrium GDP resulting from the initial \$100 billion decline in autonomous consumption. Show all steps.
    \item[(c)] Describe how the Bank of Canada could use an \textbf{open-market operation} to close this recessionary gap.
    \item[(d)] If the money multiplier is 5, how large an open-market purchase of bonds (in billions) would be required to fully offset the decline in AD? Show your reasoning.
\end{enumerate}

\vspace{1em}
\hrule
\vspace{1em}

\begin{center}
\textit{End of Practice Problem Set}
\end{center}

\end{document}
